\section{Introdução}\label{sec:intro}

\subsection{Contexto e motivação}
Sistemas operacionais e servidores expõem continuamente contadores de utilização e de atividade em baixo
nível --- carga de CPU, uso de memória, tráfego de rede, entre outros. Essa \emph{telemetria} é a base do
monitoramento: quando o comportamento de uma máquina se afasta do habitual (um vazamento de memória, um
processo descontrolado, uma falha de rede), o desvio aparece nas métricas antes de se tornar uma
interrupção visível. Detectar automaticamente esses desvios é o problema clássico de \emph{detecção de
anomalias}~\cite{chandola2009}.

A dificuldade prática é dupla. Primeiro, anomalias são raras e heterogêneas: nos dados deste trabalho
apenas 4,16\,\% dos pontos de teste são anômalos (Seção~\ref{sec:dados}), de modo que é impossível
enumerar de antemão todas as formas de falha. Segundo, o rótulo ``anômalo'' costuma vir de relatórios de
incidentes, ao passo que o comportamento \emph{normal} está disponível em abundância e sem rótulo. Isso
favorece abordagens que \emph{aprendem o que é normal} e tratam como anomalia aquilo que não se encaixa.

As redes de Hopfield~\cite{hopfield1982} são um ajuste natural a essa ideia. Trata-se de uma rede neural
recorrente totalmente conectada, cujo estado evolui, sob pesos simétricos, até um mínimo local de uma função
de energia. Os padrões gravados na matriz de pesos tornam-se atratores, e apresentar à rede uma entrada
parcial ou ruidosa a faz convergir para o padrão armazenado mais próximo~\cite{hopfieldtank1985}. Se os
padrões armazenados descrevem o funcionamento normal, esperamos que uma observação normal seja
``corrigida'' para um padrão familiar com pouco esforço, enquanto uma observação anômala permaneça distante
de qualquer atrator ou seja atraída para um estado que não corresponde a nenhuma memória. Essa diferença na
recuperação pode ser convertida em um escore de anomalia. A ideia é uma hipótese de trabalho, não um
resultado: avaliá-la é o objetivo do projeto.

\subsection{Objetivos}
\paragraph{Objetivo geral.} Avaliar a capacidade de uma rede de Hopfield de armazenar padrões de
funcionamento normal e reconhecer comportamentos que se desviam desses padrões, tratando a detecção de
anomalias em telemetria de sistemas como um problema de classificação binária (normal $\times$ anômalo).

\paragraph{Objetivos específicos.}
\begin{enumerate}
  \item Revisar o modelo de Hopfield em suas formas discreta, contínua (Hopfield--Tank) e moderna
        (Ramsauer et al.), com foco em energia, convergência, capacidade e estados espúrios.
  \item Formalizar um detector \emph{one-class} baseado em recuperação associativa, com escores definidos a
        partir de distância de Hamming e de energia.
  \item Caracterizar o conjunto de dados (SMD): amostras, atributos, rótulos, desbalanceamento e deriva.
  \item Definir o protocolo experimental: preparação, codificação binária, divisões cronológicas, métricas
        e linhas de base clássicas (Etapa~I).
  \item Definir o protocolo da Etapa~II (LLM com \emph{zero-shot} e \emph{few-shot}) sobre os mesmos dados.
  \item Formular hipóteses verificáveis e critérios para a comparação crítica dos resultados.
\end{enumerate}

\subsection{Pergunta de pesquisa}
\emph{Uma memória associativa de Hopfield, treinada somente com telemetria predominantemente normal,
separa pontos normais de pontos anômalos? E como esse desempenho se compara ao de classificadores clássicos
e ao de um LLM usado por \emph{prompting}?}

\subsection{Escopo e premissas}
\begin{notabox}{Escopo desta versão do relatório}
\begin{itemize}
  \item \textbf{Projeto e fundamentação, sem experimentos.} Este texto não contém resultados de modelos.
        As tabelas da Seção~\ref{sec:resultados} têm células vazias (\lacuna) a preencher após os
        experimentos; nenhum valor foi inventado. Os números das Seções~\ref{sec:dados} e~\ref{sec:metodo}
        são estatísticas reais do SMD (contagens, divisões e exploração dos dados) e estão identificados
        como tais.
  \item \textbf{Hopfield não consta da lista do enunciado.} Para classificação, o enunciado cita \emph{Decision
        Trees}, \emph{Random Forest}, \emph{Naive Bayes}, SVM, MLP, XGBoost e KNN~\cite{enunciado}. Adotamos
        a rede de Hopfield como \emph{modelo do tema} e \emph{Random Forest} e SVM como linhas de base
        (mínimo de 2 modelos da Etapa~I). \textbf{Confirmar este arranjo com o professor.}
  \item \textbf{``Baixo nível'' no SMD.} São 38 contadores anonimizados de utilização de recursos, sem
        chamadas de sistema nem rastros de \emph{kernel}. Coletar telemetria própria é trabalho futuro
        (Seção~\ref{sec:conclusoes}).
  \item \textbf{Etapa~II é um protocolo.} A execução exige acesso a um modelo de linguagem; aqui
        documentamos tarefa, \emph{prompt}, amostragem e métricas (Apêndice~\ref{apx:prompt}).
  \item \textbf{Propostas, não resultados.} Divisões, hiperparâmetros e linhas de base são decisões de
        projeto a validar; as derivações próprias estão sinalizadas como tais.
\end{itemize}
\end{notabox}

\subsection{Organização do texto}
A Seção~\ref{sec:teoria} apresenta a fundamentação teórica. A Seção~\ref{sec:dados} define o problema e
descreve o conjunto de dados, e a Seção~\ref{sec:metodo} detalha a metodologia, incluindo as Etapas~I e~II.
A Seção~\ref{sec:resultados} reserva as tabelas de resultados, a Seção~\ref{sec:critica} traz a comparação
crítica (hipóteses e pontos positivos e negativos) e a Seção~\ref{sec:conclusoes} fecha com conclusões e
próximos passos. Os apêndices reúnem pseudocódigo, um exemplo numérico, o \emph{prompt}, o cronograma, o
mapeamento dos requisitos do enunciado, derivações e a tabela das 28 máquinas.
